\chapter{Introducción}\label{introduccion}

En las últimas décadas, el crecimiento del parque automotor y la urbanización sostenida han convertido a la congestión vehicular en uno de los principales desafíos para la movilidad urbana. Según los reportes de flota vehicular de Argentina elaborados por la \gls{acr-afac}, y tomando 2012 como línea base, el volumen de vehículos aumentó un 21,7\% en 2017. Esta tendencia continuó durante la última década y alcanzó un incremento acumulado del 44,4\% en 2025 (véase la figura \ref{fig:flota_circulante}) \cite{AFAC2018Flota,AFAC2026Flota}.

\begin{figure}[htbp]
    \centering
    \begin{tikzpicture}[font=\sffamily\small]
    \begin{axis}[
        width=0.9\linewidth,
        height=7cm,
        xlabel={Año},
        ylabel={Flota Circulante (Millones)},
        ymin=0,
        xtick={1990, 1995, 2000, 2005, 2010, 2015, 2020, 2025},
        xticklabel style={/pgf/number format/1000 sep=},
        ymajorgrids=true,
        xmajorgrids=true,
        grid style=dashed,
        enlarge x limits=0.03
    ]
    % Área sombreada bajo la curva
    \addplot[fill=colDatoPrimario, fill opacity=0.4, draw=none] coordinates {
        (1990, 0)
        (1990, 2.203) (1992, 2.627) (1994, 3.420) (1996, 4.109)
        (1998, 4.941) (2000, 5.666) (2002, 5.943) (2004, 6.345)
        (2006, 7.081) (2008, 8.224) (2010, 9.389) (2011, 10.246)
        (2012, 10.930) (2013, 11.245) (2014, 11.520) (2015, 12.012)
        (2016, 12.503) (2017, 13.302) (2018, 13.950) (2019, 14.301)
        (2020, 14.564) (2021, 14.840) (2022, 15.079) (2023, 15.299)
        (2024, 15.552) (2025, 15.784)
        (2025, 0)
    };
    % Línea principal con marcadores
    \addplot[color=colDatoPrimario, thick, mark=*, mark options={fill=colDatoPrimario, draw=colDatoPrimario}, mark size=1.5pt] coordinates {
        (1990, 2.203) (1992, 2.627) (1994, 3.420) (1996, 4.109)
        (1998, 4.941) (2000, 5.666) (2002, 5.943) (2004, 6.345)
        (2006, 7.081) (2008, 8.224) (2010, 9.389) (2011, 10.246)
        (2012, 10.930) (2013, 11.245) (2014, 11.520) (2015, 12.012)
        (2016, 12.503) (2017, 13.302) (2018, 13.950) (2019, 14.301)
        (2020, 14.564) (2021, 14.840) (2022, 15.079) (2023, 15.299)
        (2024, 15.552) (2025, 15.784)
    };
    \end{axis}
    \end{tikzpicture}
    \caption{Evolución de la flota vehicular circulante estimada en Argentina entre 1990 y 2025, expresada en millones de vehículos. La línea y los marcadores representan los valores anuales disponibles; el área sombreada facilita la lectura de su magnitud. Elaboración propia a partir de reportes de \acrshort{acr-afac}.}\label{fig:flota_circulante}
\end{figure}

Como refleja la figura \ref{fig:flota_circulante}, la flota vehicular activa en Argentina experimentó una expansión sostenida entre 1990 y 2025, pasando de aproximadamente 2,2 millones a 15,78 millones de unidades operativas \cite{AFAC2018Flota,AFAC2026Flota}. En los reportes de \acrshort{acr-afac}, la flota circulante comprende automóviles y vehículos comerciales livianos y pesados en condiciones de circular, excluyendo acoplados y remolques; representa una estimación del parque activo bajo esa definición, no el total histórico de unidades patentadas ni la cantidad de viajes realizados. Este incremento acumulado ejerce una presión creciente sobre redes viales consolidadas que poseen una capacidad física fija, agudizando la formación de colas y demoras en las intersecciones urbanas críticas.

Las congestiones vehiculares deterioran las condiciones de operación de la red e incrementan las demoras y las detenciones de los vehículos \cite{calymayor}. En el modelado microscópico, las emisiones de gases como el monóxido de carbono (\acrshort{acr-co}), el dióxido de carbono (\acrshort{acr-co2}), los óxidos de nitrógeno (\acrshort{acr-nox}) y los hidrocarburos no combustionados (\acrshort{acr-hc}) se estiman a partir de variables cinemáticas, entre ellas la velocidad, la aceleración y el tiempo en ralentí \cite{SUMO2026_SUMO}. Este problema es especialmente relevante en ciudades con alta densidad vehicular. En este trabajo se utilizan escenarios sintéticos y un escenario urbano restringido por datos: el microcentro de la Ciudad de Mendoza, delimitado por las avenidas San Martín, Colón, Belgrano y Las Heras.

Frente a la congestión, una alternativa que no requiere modificaciones importantes de la infraestructura consiste en optimizar la gestión del tránsito mediante sistemas de control semafórico adaptativos. Los enfoques tradicionales pueden presentar limitaciones para responder a variaciones en el flujo vehicular \cite{ite_handbook}. En este contexto, los avances en \gls{acr-ia} y en particular el \gls{acr-rl} permiten modelar los semáforos como agentes que aprenden políticas de control a partir de la interacción con el entorno. Las acciones corresponden a cambios de fase y la recompensa puede definirse, por ejemplo, como el negativo del tiempo de espera acumulado. Diversos estudios han mostrado la viabilidad de este enfoque, desde algoritmos tabulares aplicados a una sola intersección \cite{Wiering2000} hasta modelos de redes coordinadas basados en redes neuronales profundas y control multiagente \cite{WeiH2019ASTSCM}.

Sin embargo, la aplicación simultánea de este enfoque en múltiples intersecciones introduce desafíos adicionales. Un esquema centralizado puede resultar poco escalable, pues el espacio de acciones conjunto crece multiplicativamente con el número de intersecciones \cite{marlbook}. En una formulación de \gls{acr-marl} con \gls{acr-il}, cada semáforo optimiza su propia intersección sin considerar explícitamente el efecto de sus decisiones sobre el resto de la red \cite{Ault2022_MARL_Limitations}. Además, como todos los agentes aprenden y modifican sus políticas al mismo tiempo, cada uno observa un entorno no estacionario. Esto dificulta la convergencia de los algoritmos de agente único \cite{foerster2017}.

Por esta razón, este trabajo aborda el problema desde el marco del \gls{acr-marl} cooperativo. Cada controlador aprende a gestionar su propia intersección, pero sus decisiones también buscan evitar que las colas y las demoras se trasladen hacia las intersecciones vecinas. En lugar de asignar una única recompensa idéntica a todos los agentes, el sistema entrega a cada uno una señal que refleja tanto la evolución de la espera local como la de su entorno inmediato. Durante el entrenamiento de las variantes centralizadas, el crítico incorpora además información conjunta de la red. De este modo, cada controlador recibe una señal vinculada con su área de influencia, sin perder de vista el funcionamiento de la red vial en su conjunto.

En este contexto, el objetivo general de este trabajo es desarrollar y evaluar un sistema de control semafórico basado en aprendizaje por refuerzo multiagente para mejorar la operación del tránsito en redes urbanas. En particular, la investigación compara el desempeño de algoritmos de entrenamiento centralizado y ejecución descentralizada (\acrshort{acr-mappo} y \acrshort{acr-happo}) con esquemas de aprendizaje independiente (\acrshort{acr-ippo} e \acrshort{acr-idqn}) bajo condiciones controladas, con el fin de evaluar estrategias de coordinación cooperativa y de tratamiento de la heterogeneidad entre intersecciones. La evaluación empírica se lleva a cabo mediante simulaciones microscópicas en escenarios sintéticos de complejidad creciente y en la red vial modelada del microcentro de la Ciudad de Mendoza, utilizando el control de tiempo fijo como ancla operacional de referencia.

Uno de los aportes metodológicos de este trabajo es representar la espera acumulada por carril mediante una transformación logarítmico-lineal acotada y discretizada en $I$ categorías con resolución no uniforme. Las categorías ofrecen mayor detalle para demoras bajas, agrupan situaciones de espera semejantes y limitan los valores extremos; así, esta representación puede facilitar el tratamiento de variaciones pequeñas y mantener acotado el rango de la observación. El uso de distintas representaciones de estado en el control semafórico con aprendizaje por refuerzo cuenta con antecedentes \cite{genders2018representations}. En esta investigación, la codificación se evalúa como parte del sistema completo; determinar su contribución específica requeriría compararla con otras codificaciones.

Además, se examina si las políticas diferenciadas son adecuadas para intersecciones con geometrías, conjuntos de acciones o demandas distintos, frente a la alternativa de compartir una política. La implementación de \acrshort{acr-happo} adapta la secuencia de actualizaciones de actores, en la que se consideran los cambios previos, con ventajas individuales y señales de recompensa locales-vecinales; en las variantes centralizadas, el crítico es compartido. Esta secuencia distingue su procedimiento de actualización del de \acrshort{acr-mappo}, sin implicar una actualización nueva del crítico ni garantizar superioridad. La comparación con políticas compartidas y con los enfoques locales de \acrshort{acr-ippo} y \acrshort{acr-idqn} permite estudiar esta cuestión en los escenarios considerados.


\section{Motivación}\label{sec:motivacion}

La motivación de este trabajo se centra en estudiar cómo coordinar controladores que aprenden en una misma red vial. El control adaptativo basado en aprendizaje por refuerzo permite ajustar decisiones a las condiciones observadas \cite{MDPI_Infrastructures2025_RLReview}. En una red de múltiples intersecciones, interesa examinar cómo esa adaptación puede incorporar los efectos de las decisiones locales sobre los cruces vecinos.

Las diferencias de demanda, geometría y fases disponibles entre intersecciones motivan, además, el estudio de políticas que puedan adaptarse a esas condiciones. La comparación de \acrshort{acr-mappo} y \acrshort{acr-happo} con los esquemas independientes \acrshort{acr-ippo} e \acrshort{acr-idqn} permite examinar distintas formas de aprendizaje y coordinación bajo una misma formulación del problema.


\section{Organización del Documento}\label{sec:organizacion_documento}

El presente trabajo se estructura de la siguiente manera:

\begin{itemize}
    \item \textbf{Capítulo 1 -- Introducción:} contextualiza la problemática de la congestión vehicular urbana, define los objetivos del proyecto y expone las motivaciones que impulsan esta investigación.
    \item \textbf{Capítulo 2 -- Marco Teórico:} presenta los conceptos fundamentales de la ingeniería de tránsito, la teoría de procesos de decisión de Markov, las redes neuronales profundas y los paradigmas de aprendizaje por refuerzo multiagente.
    \item \textbf{Capítulo 3 -- Metodología:} detalla el diseño metodológico y experimental, la calibración de la demanda vial mediante el estadístico \gls{acr-geh}, el modelado del microcentro de la Ciudad de Mendoza en \gls{acr-sumo} y la formulación algorítmica de los controladores.
    \item \textbf{Capítulo 4 -- Resultados:} presenta la evaluación empírica y comparativa de los algoritmos de control sobre redes sintéticas y el escenario urbano restringido por datos.
    \item \textbf{Capítulo 5 -- Conclusiones:} sintetiza los hallazgos principales, discute las limitaciones del modelo y propone direcciones de investigación futuras.
    \item \textbf{Apéndice A -- Arquitectura de Software:} describe las especificaciones orientadas a objetos, los diagramas de clases \gls{acr-uml} y los patrones de diseño aplicados en el marco computacional \texttt{marlTLC}.
\end{itemize}
